


\section{Méthodes}
  
\subsection{Mortalité routière}

\subsubsection{Détection de carcasses à pied et à partir d'un véhicule}

Nous avons comparé les détections de carcasses et d'individus vivants de tous les taxons \addM{vertébrés} (amphibiens, reptiles, mammifères, oiseaux) sur la chaussée de l'A10 et de la R104 en utilisant le suivi à pied par jumelles\delM{,} et le suivi par enregistrement avec caméra GoPro fixée sur un véhicule. Lors de l'échantillonnage sur l'A10, la voie Sud était séparée de la voie Nord afin de réaliser deux décomptes indépendants.Lors de ces périodes d'enregistrement, les observateurs ne notaient pas \comM[inline]{Je ne comprends pas ce que tu veux dire par les carcasses en parallèle. Clarifier.}d'observations de carcasses en parallèle pour des raisons de sécurité. Lors d'une même journée d'échantillonnage, le même segment de route était parcouru à l'aide des deux méthodes : suivi à pied par jumelles \delM{entre 15 et 30 m de la chaussée} et enregistrement par caméra GoPro fixée sur un véhicule circulant sur la chaussée. Afin de réduire la probabilité de retrait des carcasses sur la route par des prédateurs ou des automobiles, ces suivis étaient effectués dès le lever du soleil. Lors d'une journée donnée, nous avons déterminé aléatoirement l’ordre des méthodes utilisées \addM{pour un même tronçon}. Le but était de comparer l'efficacité des deux méthodes à détecter sur la route les mêmes carcasses et les mêmes individus vivants. Il était donc important de réaliser les deux suivis dans un court laps de temps (même matinée) \addM{sur le même tronçon}. L’utilisation des deux méthodes permettra de mieux estimer le nombre réel de carcasses sur la route.

Le suivi par jumelles consistait à parcourir à pied les sentiers situés à proximité des tronçons de l'A10 et de la R104, en restant entre 15 et 30 m de la chaussée. Au moins deux observateurs réalisaient l'inventaire. Le port d'un gilet de sécurité réfléchissant et d'un casque était obligatoire en tout temps pour augmenter la visibilité du personnel sur le terrain. Les jumelles utilisées étaient des Nikon Prostaff P3 8X42 (Nikon, Tokyo, Japon). Toutes les carcasses observées sur les routes étaient géolocalisées à l'aide d'un GPS GARMIN GPSMAP 79s (Garmin, Schaffhausen, Suisse). Le suivi par caméra réalisé par voiture utilisait une caméra GoPro Hero Black 11 (GoPro, San Mateo, Californie, États-Unis) afin d'enregistrer les images sur la chaussée avec une qualité 4K, un ratio 4:3 et 30 images par seconde. Lors d'une session d'échantillonnage, la caméra GoPro était insérée dans un étui imperméable transparent et fixée à un trépied magnétique (3-Footed Monster, Sydney, Australie). Le trépied magnétique était fixé au milieu du capot avant de la voiture. La caméra pouvait être contrôlée à distance par protocole Bluetooth à l'aide de l'application gratuite Quik pour Android. Nous lancions l'enregistrement lorsque le véhicule entrait dans le segment suivi jusqu'à la fin du parcours. Les fichiers vidéo ont été visionnés et les données d'observation sur la route ont été saisies en notant le moment (minutes, secondes) et la position (coordonnées géographiques) de la séquence vidéo.

Une troisième méthode de détection des carcasses, utilisant un aéronef télépiloté (ici désigné par « drone »), venait compléter les deux méthodes précédemment présentées (suivi à pied avec jumelles et caméra fixée sur véhicule). \delM{Pour rappel, c}\addM{C}ette méthode pouvait se substituer au suivi à pied lorsque les conditions météorologiques étaient favorables \addM{aux conditions de vol}. Les journées d’échantillonnage avec drone étaient moins fréquentes que celles des deux autres méthodes, car les vols nécessitent des conditions météorologiques particulières : vents inférieurs à 30 km/h et absence de précipitations. Le drone utilisé était un DJI Mavic 3 Enterprise Thermal (DJI, Shenzhen, Chine). Celui-ci a été immatriculé par Transports Canada avant les vols, et sa plaque d'immatriculation est visible sur une des pales de l'engin. Puisque le drone pesait plus de 250 g, il était nécessaire d’obtenir une certification de base pour devenir pilote. Pendant le vol, au moins une personne additionnelle aidait le pilote à conserver un contact visuel avec l’aéronef en tout temps.

Dès le lever du soleil, le drone était déployé le long de l'A10 et de la R104 une fois toutes les 2--3 semaines lors des conditions optimales de vol. Les vols ont été automatisés pour que le drone effectue des tronçons linéaires, en restant parallèle aux routes et en volant à au moins 5 m de la chaussée. Les trajets de vol étaient de 500 m chacun, le but étant de voler de part et d'autre \addM{de l'emplacement} des futurs passages, soit 250 m avant et 250 m après\delM{ les points définis de ceux-ci}. L’altitude de vol était de 40 m. Nous avons utilisé deux types de caméras, une de grand-angle et une thermique. La caméra grand-angle possède une qualité d’image de 20 mégapixels, et une lentille avec un champ de vision de 84\addM{$^{\circ}$}, un équivalent format de 24 mm, une ouverture de f/2.8 et une mise au point à 1 m. La qualité du 4K permet d’obtenir 60 images par seconde. La caméra thermique de ce drone fonctionne avec un imageur thermique, le Microbolomètre VOx non refroidi et d’un spotmètre. La lentille a un champ de vision de 61\addM{$^{\circ}$}, un équivalent format de 40 mm, une ouverture de f/1.0 et une mise au point à 5 m. La plage de mesure va de -20\addM{$^{\circ}$} à 150\addM{$^{\circ}$}C, avec une précision de +/- 2\addM{$^{\circ}$}C.

Le drone circulait à une vitesse de 2 m/s, afin de bien couvrir la chaussée lors du visionnement. La caméra du drone était orientée à 90 degrés pour faire face à la route et agrandissait l’image continuellement à un grossissement de X2. Les fichiers vidéo des vols de 500 m de qualité 4K (ratio d’image en 4:3) ont été visionnés en laboratoire et les données d'observation sur la route ont été saisies en notant le moment (minutes, secondes) de la séquence vidéo. Les carcasses ont été géolocalisées, au même titre que les deux précédentes méthodes. \delM{Le but de cette technique était de mettre en relation les données d'inventaire sur la faune retrouvée proche des futurs passages fauniques et les individus retrouvés morts sur la route.} 

\subsubsection{Persistance de carcasses sur la route}

La persistance des carcasses sur la route peut dépendre de différents facteurs tels que la consommation des carcasses par des charognards et l'intensité de circulation \addM{de véhicules sur la route}. L'équipe a combiné les suivis de mortalité à une expérience visant à évaluer la probabilité de persistance des carcasses sur la route. \delM{Le but était d’évaluer le temps passé avant qu'une carcasse ne soit retirée par des charognards.} Cette approche permettra de corriger nos estimations du nombre de carcasses sur la route afin d'évaluer le nombre réel en incluant les carcasses retirées par des charognards.

Pour ce faire, nous avons utilisé des morceaux de poulet cru \addM{de 1.5 cm x 3 cm x 3 cm} placés sur l’accotement de la chaussée. Chaque semaine, nous avons sélectionné aléatoirement deux emplacements sur l’A10 et la R104. Pour chaque semaine et chaque route, les tests ont été réalisés à deux périodes : \delM{carcasse placée} à 5h00\delM{,} \delM{et}\addM{ou} \delM{carcasse placée} à 20h00. Pour un test donné \addM{à un emplacement donné}, nous déposions un morceau de poulet \delM{de 1.5 cm x 3 cm x 3 cm} sur la chaussée du côté de l'accotement\delM{et nous prenions les coordonnées géographiques de la position du morceau}. Le morceau de poulet était congelé \delM{jusqu'à son utilisation} pour maintenir sa fraîcheur \delM{. Le}\addM{et le morceau} était dégelé quelques heures avant d'être déployé sur le terrain. \addM{À proximité du morceau de poulet, nous} avons déployé un piège photographique afin d'évaluer le temps nécessaire à un prédateur pour manger la carcasse (morceau de poulet) et aussi permettre l'identification de ce dernier. \addM{Pour ce faire, nous} avons utilisé une caméra Spypoint Force-20 (SpyPoint, Victoriaville, Québec, Canada) avec une résolution de 4 mégapixels. Le piège photographique a été installé sur un poteau en bordure de la route, à environ 1,50 m du morceau de poulet\delM{, selon l’agencement du terrain}. La caméra a été programmée avec prise de photo à intervalles réguliers de 15 secondes. Le site du test a été visité au moins 24 heures plus tard pour vérifier si la carcasse avait disparu. Le dispositif était alors démonté et déplacé ailleurs pour un autre test.

Chaque test se déroulait à un emplacement différent le long de la route pour réduire l'habituation de charognards.\comM{Quelle était la distance minimale entre deux tests de persistance pendant une même saison?} La prise de données sur le terrain consistait à noter les heures et minutes de déploiement et de retrait de caméra, les heures et minutes du moment de prédation du morceau de poulet capté par la caméra, les conditions météorologiques (pluie, température maximale) entre la pose et le retrait de la carcasse, ainsi que les coordonnées de latitude et longitude de chaque station. Ces informations permettront de mettre en relation la persistance des carcasses et les caractéristiques d'habitat à proximité de la route.

\subsubsection{Observations complémentaires}

En plus de nos propres inventaires, nous avons inclus des données d’observations de tortues sur la route provenant de science citoyenne (projet Carapace de Conservation de la nature Canada), de signalement de patrouilleurs du MTMD (SAGE), ainsi que d’informations provenant d’inventaires réalisés par des firmes de consultant avant le début du présent projet.\comM{Ces informations ont été ajoutées aux résultats?} 


\subsection{Suivi de la petite faune}

\subsubsection{Pièges photographiques}

À chaque site où un passage est prévu, l'équipe de recherche a déployé un piège photographique de chaque côté de la route, soit deux pièges par passage \comM{ajouter la figure}(FIGURE). Nous utilisons la caméra GardePro T5CF (GardePro, Hong Kong, Chine) équipée d'un capteur de mouvement passif à infrarouge très sensible et d'une résolution photo allant jusqu'à 32 mégapixels. Ces pièges photographiques ont permis de détecter la faune, notamment l'herpétofaune et les mammifères de petite et moyenne taille. Ces pièges photographiques ont été posés en mai 2024 et ont été retirés en septembre 2024.  La pose tardive du matériel en 2024 découlait de contraintes logistiques, notamment liées à la disponibilité limitée du personnel de terrain en mars et avril.

La détection des animaux s'est effectuée avec deux modes de déclenchement : la prise de photo à intervalle régulier de 5 minutes (\emph{time lapse}) et la prise d'une photo après détection de mouvement des endothermes. Afin de ne pas saturer le stockage, nous avons réglé la qualité à 4 mégapixels. Le principal avantage de ce modèle est la mise au point proche des animaux détectés avec un focus à 20 cm.

\addM{Nous avons utilisé l'outil de reconnaissance automatique MegaDetector afin de faire un tri initial des photos en deux catégories: avec animal et sans animal} \citep{norouzzadeh2021}. \addM{Cet outil repose sur l'architecture YOLO version 5 basée sur des réseaux de neurones convolutionnels et a été entraîné par Microsoft à reconnaître différents taxons. MegaDetector n'identifie pas les éléments de l'image à l'espèce, mais catégorise plutôt les éléments comme étant des animaux ou non. Pour chaque élément identifié dans une image, MegaDetector lui attribue une probabilité à être dans la bonne catégorie (i.e., animal ou non). Nous avons lancé MegaDetector sur un serveur de calcul dédié du laboratoire du professeur Mazerolle à l'Université Laval.}

\addM{Afin de réduire le nombre de faux positifs, nous avons filtré les images pour lesquelles la probabilité était $\geq 75$\%. Nous avons visionné chacune de ces images afin de valider si la détection était avérée et le cas échéant, nous avons identifié les individus sur les images au taxon le plus bas possible (famille, genre ou espèce). Ces informations ont ensuite été utilisées pour évaluer l'occurrence de différents taxons à proximité des ponceaux. De plus, ces images annotées serviront à l’entraînement d’algorithmes d’intelligence artificielle pour la reconnaissance à l'espèce, notamment à l’aide du logiciel YOLO version 11} \citep{stark2023, hopkins2024, jocher2024}.

\delM{Nous avons utilisé l'outil MegaDetector pour traiter les nombreuses images récoltées} \citep{norouzzadeh2021}. \delM{Cet outil permet de trier les images qui contiennent des photos avec animaux et de celles qui n'en contiennent pas. MegaDetector n'identifie pas à l'espèce. C'est pour cette raison que les images avec détection présumée ont été validées par un humain par la suite afin d'identifier les individus dans les images au taxon le plus bas possible (famille, genre ou espèce).}


\subsubsection{Enregistrements accoustiques}

En 2023 et 2024, l'équipe a installé des enregistreurs acoustiques automatisés à la bordure de milieux humides se trouvant dans l'emprise du projet afin de détecter les anoures et d'évaluer leur degré \addM{d'activité} \comM{Ajouter figure d'un enregistreur sur le terrain}. L'équipe les a placés dès le début de la période de reproduction des anoures fin mars 2023 et 2024, et les a retirés fin septembre 2023 et fin août 2024 afin de couvrir \delM{toutes} les périodes d'activités des différentes espèces. Les appareils sont des AudioMoths \citep{hill2019} de 7 cm X 8 cm X 3 cm. Ils sont placés à environ 1,5 m du sol et fixés à des arbres \addM{ou des poteaux}. Les appareils sont programmés pour enregistrer les sons durant 3 minutes, à quatre reprises dans la journée (20h00, 00h00, 10h00, 14h00). Ces plages sont celles pendant lesquelles différentes espèces d’anoures sont actives, incluant les rainettes faux-grillon \citep{fayard2022}. Les plages d'enregistrement ont été analysées à l’aide de BirdNET qui est un outil de reconnaissance utilisant l’intelligence artificielle \citep{kahl2021}. Cet outil permet d’identifier sur chaque enregistrement sonore les espèces d’anoures présentes. Toutefois, BirdNET est entraîné pour reconnaître huit des dix espèces d’anoures du Québec pour le moment. La grenouille léopard (\emph{Lithobates pipiens}) et la grenouille du Nord (\emph{Lithobates septentrionalis}) ne figurent pas sur la liste des espèces reconnues par ce modèle global. \addM{Afin de} détecter ces deux espèces, nous avons entraîné à l'automne 2025 un classificateur personnalisé à partir de BirdNET. \addM{La méthode d'entraînement de BirdNET est détaillée à la partie II du rapport}.

\subsubsection{Recherche par drones}

Nous avons réalisé des suivis par drones au-dessus des zones humides de l'A10 et de la R104, afin de détecter les tortues qui thermorégulent dans les milieux humides \citep{aota2021, bogolin2021, bevan2018}. Nous avons utilisé le même drone décrit plus haut pour ces inventaires\delM{, en suivant le même protocole pour les observateurs visuels}\comM{Pas clair de quoi il est question avec les observateurs visuels. Tu veux dire que le drone était toujours dans le champ de vision d'un observateur? Clarifier.}. Cette technologie permet d'éviter l'effarouchement des tortues, en respectant les recommandations de vol à une altitude minimale de 20 m au-dessus des milieux humides \citep{charbonneau2021}. Les survols en drone permettent aussi une couverture complète de la zone d’inventaire, et facilitent l’accès à des secteurs qui auraient été inaccessibles à pied.

Transports Canada réglemente les zones où le vol de drone est autorisé. La zone de vol doit se situer à plus de 9 km de tout aéroport, héliport ou aérodrome. De même, les vols sont interdits au-dessus d’espaces réglementés, comme une base militaire, une prison ou une aire protégée comme celle du Bois de Brossard. Les secteurs d’inventaire que nous avons sélectionnés n’ont pas ce type de contrainte. Nous avons cherché des aires de décollage et d’atterrissage sécuritaires. Ces suivis ont été réalisés à différentes périodes de la saison d'activité des tortues entre mai et octobre 2024 \citep{boyer1965, obbard1981, lefevre1995}.

Nous avons suivi la méthodologie de recherche active proposée dans le protocole standardisé de détection et d'identification des tortues d'eau douce à l'aide de drones au Québec \citep{mffp2021}. Nous avons rempli le "formulaire de prise de données" figurant dans l'annexe C de ce même protocole après chaque sortie \citep{mffp2021}. Les sorties se faisaient selon les conditions météorologiques suivantes : journée ensoleillée et couverture nuageuse pouvant aller jusqu’à 75 \%, température de l’air entre 10 et 25{\textdegree}C, journée sans vent ou jusqu’à une légère brise (niveau 2 sur l’échelle de Beaufort). La caméra thermique du drone était utilisée pour \delM{connaître}\addM{mesurer} la température des tortues observées, \addM{de leur substrat, ainsi que de l’eau dans laquelle} \delM{u milieu où} elles se trouvaient. Pendant une mission de vol, les milieux humides ont été survolés à basse vitesse à une altitude variant de 30 à 40 m. Lorsqu'une forme ressemblant à une tortue était repérée, nous avons utilisé le zoom du drone allant jusqu'à un grossissement de 56x pour confirmer qu'il s’agissait d'une tortue. En plus de la recherche active, lors d’une même visite nous avons lancé des vols automatisés qui filmaient l’intégralité du milieu humide en suivant toujours le même chemin à une altitude fixe de 40 m et un angle de caméra fixé à 90\textdegree. Le but des vols automatisés était de constituer une banque de vidéos qui pourront servir à entraîner \delM{une}\addM{un outil d'}intelligence artificielle (IA) pour identifier les tortues de manière automatique. \delM{Étant donné que la recherche active est plus longue que les vols automatisés,  nous voudrions utiliser l’IA sur les vols automatisés de tous les sites suivis.}

Concernant l'A10, nous avons priorisé les zones où des tortues avaient été vues lors d'inventaires par drones en 2022 \addM{par une autre équipe}\delM{pour ce même projet} \citep{snc2023}. Des polygones avaient été choisis en fonction de la présence d'étangs ou de marécages visibles par l'imagerie satellitaire. D’autres milieux humides ont été survolés lorsqu’ils semblaient favorables à l’habitat des tortues. Du côté de la R104, nous avons prospecté le long de la rivière Saint-Jacques, où il existe des données historiques de tortues peintes et de tortues serpentines qui traversent la R104 (\citealt{fayard2024a}).

\subsubsection{Stations de plaques à reptiles}
Nous avons déployé des grilles de refuges artificiels à moins de 30 mètres de la route afin d'inventorier différentes espèces de couleuvres \citep{dodd2016}. À chaque site où un passage était prévu, l'équipe a installé une station de chaque côté de la route. Une station consistait en quatre panneaux de contreplaqué\delM{s} d'épinette non traité\delM{s} de 60 cm x 60 cm x 2 cm posés sur un carré de géotextile noir de 73 cm x 73 cm. Afin de connaître également les espèces qui fréquentaient les habitats qui n’étaient pas adjacents aux passages fauniques prévus, nous avons ajouté des stations qui serviraient de témoins entre chaque passage. \addM{Un total de 16 stations de plaques à reptiles ont été déployées sur}\delM{Sur} l'A10\delM{, cela faisait un total de 16 stations de plaques à reptiles}, soit huit stations de chaque côté de la route. Pour la R104, il y avait un total de huit stations, soit quatre de chaque côté de la route (cf. Carte 4, \citealt{fayard2024a}).

En nous basant sur les recommandations du \og Protocole standardisé pour les inventaires de couleuvres \fg{} du MELCCFP, nous avons placé les grilles d'inventaires aux endroits propices aux couleuvres, i.e., dans des habitats dégagés, exposés au moins 8 heures par jour au soleil\comM[inline]{Ça prend la citation du rapport en question.}. Chaque semaine, l'équipe vérifiait les stations dans tous les sites en inspectant et soulevant les panneaux de bois et le géotextile. Nous avons privilégié les visites lorsque les températures étaient entre 15 et 25{\textdegree}C. Au-delà de 25{\textdegree}C, la température était trop chaude et les couleuvres n'utilisaient plus les abris.

\subsubsection{Autres observations fauniques}

Des observations additionnelles de l’herpétofaune, des mammifères et des oiseaux ont été collectées de façon opportuniste lors des visites sur le terrain. Une photo avec un cellulaire intelligent a été prise lorsque c'était possible, permettant de géoréférencer le point d’observation avec les métadonnées des photos.
